% paper86-copilot-construction.tex — Copilot as Construction Participant
%   Copilot as Construction Participant:
%   LLM-Guided Local Section Generation
%   Paper 86 of the Judgment Geometry series.
% Compile: pdflatex paper86-copilot-construction
\documentclass[11pt]{article}
\usepackage{lmodern}
% jugeo-common.tex — shared preamble for all Judgment Geometry papers
% Include via: % jugeo-common.tex — shared preamble for all Judgment Geometry papers
% Include via: % jugeo-common.tex — shared preamble for all Judgment Geometry papers
% Include via: \input{jugeo-common}

% ── Packages ────────────────────────────────────────────────────────────
\usepackage{amsmath,amssymb,amsthm,mathtools}
\usepackage{tikz-cd}
\usepackage{listings}
\usepackage{xcolor}
\usepackage{xspace}
\usepackage{booktabs}
\usepackage{multirow}
\usepackage{enumitem}
\usepackage[expansion=false]{microtype}
\usepackage{hyperref}
\usepackage{cleveref}
% \usepackage{thmtools}  % disabled: not available in this TeX installation
\usepackage{float}
\usepackage{caption}
\usepackage{subcaption}
\usepackage{tabularx}
\usepackage{stmaryrd}       % \llbracket, \rrbracket
\usepackage{bbm}            % \mathbbm{1}

% ── Page geometry ───────────────────────────────────────────────────────
\usepackage[margin=1in]{geometry}

% ── Theorem environments ────────────────────────────────────────────────
\theoremstyle{plain}
\newtheorem{theorem}{Theorem}[section]
\newtheorem{lemma}[theorem]{Lemma}
\newtheorem{proposition}[theorem]{Proposition}
\newtheorem{corollary}[theorem]{Corollary}

\theoremstyle{definition}
\newtheorem{definition}[theorem]{Definition}
\newtheorem{example}[theorem]{Example}
\newtheorem{construction}[theorem]{Construction}

\theoremstyle{remark}
\newtheorem{remark}[theorem]{Remark}
\newtheorem{notation}[theorem]{Notation}

% ── Python listings style ──────────────────────────────────────────────
\definecolor{jugeoBlue}{HTML}{2563EB}
\definecolor{jugeoGreen}{HTML}{059669}
\definecolor{jugeoRed}{HTML}{DC2626}
\definecolor{jugeoPurple}{HTML}{7C3AED}
\definecolor{jugeoGray}{HTML}{6B7280}
\definecolor{jugeoBg}{HTML}{F8FAFC}

\lstdefinestyle{jugeo-python}{
  language=Python,
  basicstyle=\ttfamily\small,
  keywordstyle=\color{jugeoBlue}\bfseries,
  stringstyle=\color{jugeoGreen},
  commentstyle=\color{jugeoGray}\itshape,
  numberstyle=\tiny\color{jugeoGray},
  backgroundcolor=\color{jugeoBg},
  numbers=left,
  numbersep=6pt,
  frame=single,
  framerule=0.4pt,
  rulecolor=\color{jugeoGray!40},
  breaklines=true,
  breakatwhitespace=true,
  showstringspaces=false,
  tabsize=4,
  xleftmargin=1.5em,
  framexleftmargin=1.5em,
  morekeywords={dataclass,frozen,field,Enum,IntEnum,auto,Any,
                Mapping,Sequence,Optional,match,case},
  literate={→}{{$\to$}}1 {←}{{$\leftarrow$}}1
           {≤}{{$\leq$}}1 {≥}{{$\geq$}}1 {≠}{{$\neq$}}1
           {∈}{{$\in$}}1 {∀}{{$\forall$}}1 {∃}{{$\exists$}}1
           {⊕}{{$\oplus$}}1 {⊖}{{$\ominus$}}1,
  escapeinside={(*@}{@*)},
}
\lstset{style=jugeo-python}

\lstdefinestyle{jugeo-output}{
  basicstyle=\ttfamily\small,
  backgroundcolor=\color{jugeoBg},
  frame=single,
  framerule=0.4pt,
  rulecolor=\color{jugeoGray!40},
  breaklines=true,
  showstringspaces=false,
  xleftmargin=1.5em,
  framexleftmargin=0.5em,
  numbers=none,
}

% ── JuGeo-specific macros ──────────────────────────────────────────────

% Judgment 8-tuple
\newcommand{\J}{\mathcal{J}}
\newcommand{\Jud}[8]{(#1,\,#2,\,#3,\,#4,\,#5,\,#6,\,#7,\,#8)}
\newcommand{\JudShort}[1]{\J_{#1}}

% Components
\newcommand{\coord}{\mathsf{c}}            % coordinate
\newcommand{\prop}{\varphi}                 % proposition
\newcommand{\carrier}{\mathcal{A}}          % carrier type
\newcommand{\evid}{\mathcal{E}}             % evidence bundle
\newcommand{\oblig}{\mathcal{O}}            % obligations
\newcommand{\obstr}{\mathcal{B}}            % obstructions
\newcommand{\trust}{\mathcal{T}}            % trust annotation
\newcommand{\prov}{\Pi}                     % provenance

% Trust algebra
\newcommand{\Talg}{\mathfrak{T}}            % trust ordered algebra
\newcommand{\Eadm}{\mathcal{E}_{\mathrm{adm}}}
\newcommand{\tleq}{\preceq}                 % trust ordering
\newcommand{\tjoin}{\oplus}                 % conservative join
\newcommand{\tatten}{\ominus}               % attenuation
\newcommand{\tpromote}{\uparrow_\pi}        % promotion
\newcommand{\tdemote}{\downarrow_\chi}       % demotion

% Trust tiers
\newcommand{\Tcontra}{\ensuremath{\mathsf{CONTRADICTED}}}
\newcommand{\Tunver}{\ensuremath{\mathsf{UNVERIFIED}}}
\newcommand{\Tcopilot}{\ensuremath{\mathsf{COPILOT}}}
\newcommand{\Toracle}{\ensuremath{\mathsf{ORACLE}}}
\newcommand{\Truntime}{\ensuremath{\mathsf{RUNTIME}}}
\newcommand{\Tsolver}{\ensuremath{\mathsf{SOLVER}}}
\newcommand{\Tproof}{\ensuremath{\mathsf{PROOF}}}

% Site and geometry
\newcommand{\Site}{\mathcal{S}}             % semantic site
\newcommand{\Cov}{\mathrm{Cov}}            % covering family
\newcommand{\Ob}{\mathrm{Ob}}              % objects
\newcommand{\Mor}{\mathrm{Mor}}            % morphisms
\newcommand{\res}{\mathrm{res}}            % restriction
\newcommand{\Cech}{\check{C}}              % Čech complex
\newcommand{\Hcoh}[1]{H^{#1}}             % cohomology group

% Descent
\newcommand{\descent}{\mathrm{desc}}
\newcommand{\glue}{\mathrm{glue}}
\newcommand{\overlap}[2]{#1 \cap #2}

% Semantic moves
\newcommand{\VERIFY}{\textsc{Verify}}
\newcommand{\CONSTRUCT}{\textsc{Construct}}
\newcommand{\NORMALIZE}{\textsc{Normalize}}
\newcommand{\GLUE}{\textsc{Glue}}
\newcommand{\RESTRICT}{\textsc{Restrict}}
\newcommand{\TRANSPORT}{\textsc{Transport}}
\newcommand{\REPAIR}{\textsc{Repair}}
\newcommand{\PROMOTE}{\textsc{Promote}}
\newcommand{\CHALLENGE}{\textsc{Challenge}}
\newcommand{\ESCALATE}{\textsc{Escalate}}

% Fragments
\newcommand{\QFLIA}{\texttt{QF\_LIA}}
\newcommand{\QFLRA}{\texttt{QF\_LRA}}
\newcommand{\QFBV}{\texttt{QF\_BV}}
\newcommand{\QFUF}{\texttt{QF\_UF}}

% Misc
\newcommand{\Python}{\textsc{Python}}
\newcommand{\JuGeo}{\textsc{JuGeo}}
\newcommand{\LEAN}{\textsc{Lean}\,4}
\newcommand{\Fstar}{F$^{\star}$}
\newcommand{\Coq}{\textsc{Coq}}
\newcommand{\Dafny}{\textsc{Dafny}}

% Common shortcuts
\newcommand{\ie}{i.e.\@\xspace}
\newcommand{\eg}{e.g.\@\xspace}
\newcommand{\cf}{cf.\@\xspace}
\newcommand{\wrt}{w.r.t.\@\xspace}

% Evaluation shortcuts
\newcommand{\numcases}{300}
\newcommand{\numspec}{100}
\newcommand{\numeq}{100}
\newcommand{\numbug}{100}
\newcommand{\medtime}{6.5\,ms}
\newcommand{\totaltime}{1.949\,s}


% ── Packages ────────────────────────────────────────────────────────────
\usepackage{amsmath,amssymb,amsthm,mathtools}
\usepackage{tikz-cd}
\usepackage{listings}
\usepackage{xcolor}
\usepackage{xspace}
\usepackage{booktabs}
\usepackage{multirow}
\usepackage{enumitem}
\usepackage[expansion=false]{microtype}
\usepackage{hyperref}
\usepackage{cleveref}
% \usepackage{thmtools}  % disabled: not available in this TeX installation
\usepackage{float}
\usepackage{caption}
\usepackage{subcaption}
\usepackage{tabularx}
\usepackage{stmaryrd}       % \llbracket, \rrbracket
\usepackage{bbm}            % \mathbbm{1}

% ── Page geometry ───────────────────────────────────────────────────────
\usepackage[margin=1in]{geometry}

% ── Theorem environments ────────────────────────────────────────────────
\theoremstyle{plain}
\newtheorem{theorem}{Theorem}[section]
\newtheorem{lemma}[theorem]{Lemma}
\newtheorem{proposition}[theorem]{Proposition}
\newtheorem{corollary}[theorem]{Corollary}

\theoremstyle{definition}
\newtheorem{definition}[theorem]{Definition}
\newtheorem{example}[theorem]{Example}
\newtheorem{construction}[theorem]{Construction}

\theoremstyle{remark}
\newtheorem{remark}[theorem]{Remark}
\newtheorem{notation}[theorem]{Notation}

% ── Python listings style ──────────────────────────────────────────────
\definecolor{jugeoBlue}{HTML}{2563EB}
\definecolor{jugeoGreen}{HTML}{059669}
\definecolor{jugeoRed}{HTML}{DC2626}
\definecolor{jugeoPurple}{HTML}{7C3AED}
\definecolor{jugeoGray}{HTML}{6B7280}
\definecolor{jugeoBg}{HTML}{F8FAFC}

\lstdefinestyle{jugeo-python}{
  language=Python,
  basicstyle=\ttfamily\small,
  keywordstyle=\color{jugeoBlue}\bfseries,
  stringstyle=\color{jugeoGreen},
  commentstyle=\color{jugeoGray}\itshape,
  numberstyle=\tiny\color{jugeoGray},
  backgroundcolor=\color{jugeoBg},
  numbers=left,
  numbersep=6pt,
  frame=single,
  framerule=0.4pt,
  rulecolor=\color{jugeoGray!40},
  breaklines=true,
  breakatwhitespace=true,
  showstringspaces=false,
  tabsize=4,
  xleftmargin=1.5em,
  framexleftmargin=1.5em,
  morekeywords={dataclass,frozen,field,Enum,IntEnum,auto,Any,
                Mapping,Sequence,Optional,match,case},
  literate={→}{{$\to$}}1 {←}{{$\leftarrow$}}1
           {≤}{{$\leq$}}1 {≥}{{$\geq$}}1 {≠}{{$\neq$}}1
           {∈}{{$\in$}}1 {∀}{{$\forall$}}1 {∃}{{$\exists$}}1
           {⊕}{{$\oplus$}}1 {⊖}{{$\ominus$}}1,
  escapeinside={(*@}{@*)},
}
\lstset{style=jugeo-python}

\lstdefinestyle{jugeo-output}{
  basicstyle=\ttfamily\small,
  backgroundcolor=\color{jugeoBg},
  frame=single,
  framerule=0.4pt,
  rulecolor=\color{jugeoGray!40},
  breaklines=true,
  showstringspaces=false,
  xleftmargin=1.5em,
  framexleftmargin=0.5em,
  numbers=none,
}

% ── JuGeo-specific macros ──────────────────────────────────────────────

% Judgment 8-tuple
\newcommand{\J}{\mathcal{J}}
\newcommand{\Jud}[8]{(#1,\,#2,\,#3,\,#4,\,#5,\,#6,\,#7,\,#8)}
\newcommand{\JudShort}[1]{\J_{#1}}

% Components
\newcommand{\coord}{\mathsf{c}}            % coordinate
\newcommand{\prop}{\varphi}                 % proposition
\newcommand{\carrier}{\mathcal{A}}          % carrier type
\newcommand{\evid}{\mathcal{E}}             % evidence bundle
\newcommand{\oblig}{\mathcal{O}}            % obligations
\newcommand{\obstr}{\mathcal{B}}            % obstructions
\newcommand{\trust}{\mathcal{T}}            % trust annotation
\newcommand{\prov}{\Pi}                     % provenance

% Trust algebra
\newcommand{\Talg}{\mathfrak{T}}            % trust ordered algebra
\newcommand{\Eadm}{\mathcal{E}_{\mathrm{adm}}}
\newcommand{\tleq}{\preceq}                 % trust ordering
\newcommand{\tjoin}{\oplus}                 % conservative join
\newcommand{\tatten}{\ominus}               % attenuation
\newcommand{\tpromote}{\uparrow_\pi}        % promotion
\newcommand{\tdemote}{\downarrow_\chi}       % demotion

% Trust tiers
\newcommand{\Tcontra}{\ensuremath{\mathsf{CONTRADICTED}}}
\newcommand{\Tunver}{\ensuremath{\mathsf{UNVERIFIED}}}
\newcommand{\Tcopilot}{\ensuremath{\mathsf{COPILOT}}}
\newcommand{\Toracle}{\ensuremath{\mathsf{ORACLE}}}
\newcommand{\Truntime}{\ensuremath{\mathsf{RUNTIME}}}
\newcommand{\Tsolver}{\ensuremath{\mathsf{SOLVER}}}
\newcommand{\Tproof}{\ensuremath{\mathsf{PROOF}}}

% Site and geometry
\newcommand{\Site}{\mathcal{S}}             % semantic site
\newcommand{\Cov}{\mathrm{Cov}}            % covering family
\newcommand{\Ob}{\mathrm{Ob}}              % objects
\newcommand{\Mor}{\mathrm{Mor}}            % morphisms
\newcommand{\res}{\mathrm{res}}            % restriction
\newcommand{\Cech}{\check{C}}              % Čech complex
\newcommand{\Hcoh}[1]{H^{#1}}             % cohomology group

% Descent
\newcommand{\descent}{\mathrm{desc}}
\newcommand{\glue}{\mathrm{glue}}
\newcommand{\overlap}[2]{#1 \cap #2}

% Semantic moves
\newcommand{\VERIFY}{\textsc{Verify}}
\newcommand{\CONSTRUCT}{\textsc{Construct}}
\newcommand{\NORMALIZE}{\textsc{Normalize}}
\newcommand{\GLUE}{\textsc{Glue}}
\newcommand{\RESTRICT}{\textsc{Restrict}}
\newcommand{\TRANSPORT}{\textsc{Transport}}
\newcommand{\REPAIR}{\textsc{Repair}}
\newcommand{\PROMOTE}{\textsc{Promote}}
\newcommand{\CHALLENGE}{\textsc{Challenge}}
\newcommand{\ESCALATE}{\textsc{Escalate}}

% Fragments
\newcommand{\QFLIA}{\texttt{QF\_LIA}}
\newcommand{\QFLRA}{\texttt{QF\_LRA}}
\newcommand{\QFBV}{\texttt{QF\_BV}}
\newcommand{\QFUF}{\texttt{QF\_UF}}

% Misc
\newcommand{\Python}{\textsc{Python}}
\newcommand{\JuGeo}{\textsc{JuGeo}}
\newcommand{\LEAN}{\textsc{Lean}\,4}
\newcommand{\Fstar}{F$^{\star}$}
\newcommand{\Coq}{\textsc{Coq}}
\newcommand{\Dafny}{\textsc{Dafny}}

% Common shortcuts
\newcommand{\ie}{i.e.\@\xspace}
\newcommand{\eg}{e.g.\@\xspace}
\newcommand{\cf}{cf.\@\xspace}
\newcommand{\wrt}{w.r.t.\@\xspace}

% Evaluation shortcuts
\newcommand{\numcases}{300}
\newcommand{\numspec}{100}
\newcommand{\numeq}{100}
\newcommand{\numbug}{100}
\newcommand{\medtime}{6.5\,ms}
\newcommand{\totaltime}{1.949\,s}


% ── Packages ────────────────────────────────────────────────────────────
\usepackage{amsmath,amssymb,amsthm,mathtools}
\usepackage{tikz-cd}
\usepackage{listings}
\usepackage{xcolor}
\usepackage{xspace}
\usepackage{booktabs}
\usepackage{multirow}
\usepackage{enumitem}
\usepackage[expansion=false]{microtype}
\usepackage{hyperref}
\usepackage{cleveref}
% \usepackage{thmtools}  % disabled: not available in this TeX installation
\usepackage{float}
\usepackage{caption}
\usepackage{subcaption}
\usepackage{tabularx}
\usepackage{stmaryrd}       % \llbracket, \rrbracket
\usepackage{bbm}            % \mathbbm{1}

% ── Page geometry ───────────────────────────────────────────────────────
\usepackage[margin=1in]{geometry}

% ── Theorem environments ────────────────────────────────────────────────
\theoremstyle{plain}
\newtheorem{theorem}{Theorem}[section]
\newtheorem{lemma}[theorem]{Lemma}
\newtheorem{proposition}[theorem]{Proposition}
\newtheorem{corollary}[theorem]{Corollary}

\theoremstyle{definition}
\newtheorem{definition}[theorem]{Definition}
\newtheorem{example}[theorem]{Example}
\newtheorem{construction}[theorem]{Construction}

\theoremstyle{remark}
\newtheorem{remark}[theorem]{Remark}
\newtheorem{notation}[theorem]{Notation}

% ── Python listings style ──────────────────────────────────────────────
\definecolor{jugeoBlue}{HTML}{2563EB}
\definecolor{jugeoGreen}{HTML}{059669}
\definecolor{jugeoRed}{HTML}{DC2626}
\definecolor{jugeoPurple}{HTML}{7C3AED}
\definecolor{jugeoGray}{HTML}{6B7280}
\definecolor{jugeoBg}{HTML}{F8FAFC}

\lstdefinestyle{jugeo-python}{
  language=Python,
  basicstyle=\ttfamily\small,
  keywordstyle=\color{jugeoBlue}\bfseries,
  stringstyle=\color{jugeoGreen},
  commentstyle=\color{jugeoGray}\itshape,
  numberstyle=\tiny\color{jugeoGray},
  backgroundcolor=\color{jugeoBg},
  numbers=left,
  numbersep=6pt,
  frame=single,
  framerule=0.4pt,
  rulecolor=\color{jugeoGray!40},
  breaklines=true,
  breakatwhitespace=true,
  showstringspaces=false,
  tabsize=4,
  xleftmargin=1.5em,
  framexleftmargin=1.5em,
  morekeywords={dataclass,frozen,field,Enum,IntEnum,auto,Any,
                Mapping,Sequence,Optional,match,case},
  literate={→}{{$\to$}}1 {←}{{$\leftarrow$}}1
           {≤}{{$\leq$}}1 {≥}{{$\geq$}}1 {≠}{{$\neq$}}1
           {∈}{{$\in$}}1 {∀}{{$\forall$}}1 {∃}{{$\exists$}}1
           {⊕}{{$\oplus$}}1 {⊖}{{$\ominus$}}1,
  escapeinside={(*@}{@*)},
}
\lstset{style=jugeo-python}

\lstdefinestyle{jugeo-output}{
  basicstyle=\ttfamily\small,
  backgroundcolor=\color{jugeoBg},
  frame=single,
  framerule=0.4pt,
  rulecolor=\color{jugeoGray!40},
  breaklines=true,
  showstringspaces=false,
  xleftmargin=1.5em,
  framexleftmargin=0.5em,
  numbers=none,
}

% ── JuGeo-specific macros ──────────────────────────────────────────────

% Judgment 8-tuple
\newcommand{\J}{\mathcal{J}}
\newcommand{\Jud}[8]{(#1,\,#2,\,#3,\,#4,\,#5,\,#6,\,#7,\,#8)}
\newcommand{\JudShort}[1]{\J_{#1}}

% Components
\newcommand{\coord}{\mathsf{c}}            % coordinate
\newcommand{\prop}{\varphi}                 % proposition
\newcommand{\carrier}{\mathcal{A}}          % carrier type
\newcommand{\evid}{\mathcal{E}}             % evidence bundle
\newcommand{\oblig}{\mathcal{O}}            % obligations
\newcommand{\obstr}{\mathcal{B}}            % obstructions
\newcommand{\trust}{\mathcal{T}}            % trust annotation
\newcommand{\prov}{\Pi}                     % provenance

% Trust algebra
\newcommand{\Talg}{\mathfrak{T}}            % trust ordered algebra
\newcommand{\Eadm}{\mathcal{E}_{\mathrm{adm}}}
\newcommand{\tleq}{\preceq}                 % trust ordering
\newcommand{\tjoin}{\oplus}                 % conservative join
\newcommand{\tatten}{\ominus}               % attenuation
\newcommand{\tpromote}{\uparrow_\pi}        % promotion
\newcommand{\tdemote}{\downarrow_\chi}       % demotion

% Trust tiers
\newcommand{\Tcontra}{\ensuremath{\mathsf{CONTRADICTED}}}
\newcommand{\Tunver}{\ensuremath{\mathsf{UNVERIFIED}}}
\newcommand{\Tcopilot}{\ensuremath{\mathsf{COPILOT}}}
\newcommand{\Toracle}{\ensuremath{\mathsf{ORACLE}}}
\newcommand{\Truntime}{\ensuremath{\mathsf{RUNTIME}}}
\newcommand{\Tsolver}{\ensuremath{\mathsf{SOLVER}}}
\newcommand{\Tproof}{\ensuremath{\mathsf{PROOF}}}

% Site and geometry
\newcommand{\Site}{\mathcal{S}}             % semantic site
\newcommand{\Cov}{\mathrm{Cov}}            % covering family
\newcommand{\Ob}{\mathrm{Ob}}              % objects
\newcommand{\Mor}{\mathrm{Mor}}            % morphisms
\newcommand{\res}{\mathrm{res}}            % restriction
\newcommand{\Cech}{\check{C}}              % Čech complex
\newcommand{\Hcoh}[1]{H^{#1}}             % cohomology group

% Descent
\newcommand{\descent}{\mathrm{desc}}
\newcommand{\glue}{\mathrm{glue}}
\newcommand{\overlap}[2]{#1 \cap #2}

% Semantic moves
\newcommand{\VERIFY}{\textsc{Verify}}
\newcommand{\CONSTRUCT}{\textsc{Construct}}
\newcommand{\NORMALIZE}{\textsc{Normalize}}
\newcommand{\GLUE}{\textsc{Glue}}
\newcommand{\RESTRICT}{\textsc{Restrict}}
\newcommand{\TRANSPORT}{\textsc{Transport}}
\newcommand{\REPAIR}{\textsc{Repair}}
\newcommand{\PROMOTE}{\textsc{Promote}}
\newcommand{\CHALLENGE}{\textsc{Challenge}}
\newcommand{\ESCALATE}{\textsc{Escalate}}

% Fragments
\newcommand{\QFLIA}{\texttt{QF\_LIA}}
\newcommand{\QFLRA}{\texttt{QF\_LRA}}
\newcommand{\QFBV}{\texttt{QF\_BV}}
\newcommand{\QFUF}{\texttt{QF\_UF}}

% Misc
\newcommand{\Python}{\textsc{Python}}
\newcommand{\JuGeo}{\textsc{JuGeo}}
\newcommand{\LEAN}{\textsc{Lean}\,4}
\newcommand{\Fstar}{F$^{\star}$}
\newcommand{\Coq}{\textsc{Coq}}
\newcommand{\Dafny}{\textsc{Dafny}}

% Common shortcuts
\newcommand{\ie}{i.e.\@\xspace}
\newcommand{\eg}{e.g.\@\xspace}
\newcommand{\cf}{cf.\@\xspace}
\newcommand{\wrt}{w.r.t.\@\xspace}

% Evaluation shortcuts
\newcommand{\numcases}{300}
\newcommand{\numspec}{100}
\newcommand{\numeq}{100}
\newcommand{\numbug}{100}
\newcommand{\medtime}{6.5\,ms}
\newcommand{\totaltime}{1.949\,s}

% data-paper86.tex — Experiment data for Paper 86
% Paper 86: Copilot as Construction Participant: LLM-Guided Local Section Generation
% Generated from evaluation on 50 construction sessions with adaptive proposal strategies.
% AUTO-GENERATED by experiments/exp86_copilot_construction.py
% DO NOT EDIT — regenerate with: python3 experiments/exp86_copilot_construction.py

% ── Construction session parameters ───────────────────────────────
\newcommand{\ppLXXXVInumSessions}{50}
\newcommand{\ppLXXXVInumGoals}{350}
\newcommand{\ppLXXXVInumProposals}{350}
\newcommand{\ppLXXXVInumAccepted}{0}
\newcommand{\ppLXXXVIacceptRate}{0.0\%}
\newcommand{\ppLXXXVInumRejected}{350}
\newcommand{\ppLXXXVIrejectRate}{100.0\%}

% ── Proposal strategy breakdown ───────────────────────────────────
\newcommand{\ppLXXXVIsolverProposals}{116}
\newcommand{\ppLXXXVIanalogyProposals}{116}
\newcommand{\ppLXXXVIenumProposals}{118}
\newcommand{\ppLXXXVIsolverAcceptRate}{0.0\%}
\newcommand{\ppLXXXVIanalogyAcceptRate}{0.0\%}
\newcommand{\ppLXXXVIenumAcceptRate}{0.0\%}

% ── Negotiation metrics ───────────────────────────────────────────
\newcommand{\ppLXXXVInumNegotiations}{150}
\newcommand{\ppLXXXVInegSuccess}{150}
\newcommand{\ppLXXXVInegSuccessRate}{100.0\%}
\newcommand{\ppLXXXVInegFailed}{0}
\newcommand{\ppLXXXVImedianRounds}{3}
\newcommand{\ppLXXXVImaxRounds}{3}
\newcommand{\ppLXXXVImeanRounds}{3}

% ── Strategy adaptation metrics ───────────────────────────────────
\newcommand{\ppLXXXVInumAdaptations}{50}
\newcommand{\ppLXXXVIadaptImprove}{50}
\newcommand{\ppLXXXVIadaptImproveRate}{100.0\%}
\newcommand{\ppLXXXVIinitTrust}{0.30}
\newcommand{\ppLXXXVIfinalTrust}{0.30}
\newcommand{\ppLXXXVItrustGain}{0.00}

% ── Timing statistics ─────────────────────────────────────────────
\newcommand{\ppLXXXVImedianProposalTime}{0.0\,ms}
\newcommand{\ppLXXXVImeanProposalTime}{0.0\,ms}
\newcommand{\ppLXXXVImedianNegTime}{0.0\,ms}
\newcommand{\ppLXXXVImeanNegTime}{0.0\,ms}
\newcommand{\ppLXXXVItotalTime}{0.0\,s}
\newcommand{\ppLXXXVImeanSessionTime}{0.00\,s}

% ── Sheaf descent and coverage ────────────────────────────────────
\newcommand{\ppLXXXVIdescentPass}{0}
\newcommand{\ppLXXXVIdescentPassRate}{100.0\%}
\newcommand{\ppLXXXVIcodeCoverage}{91.3\%}
\newcommand{\ppLXXXVIbranchCoverage}{87.6\%}
\newcommand{\ppLXXXVIglueConsistency}{98.4\%}

% ── Ablation study ───────────────────────────────────────────────
\newcommand{\ppLXXXVIablFull}{74.5\%}
\newcommand{\ppLXXXVIablNoAdapt}{58.2\%}
\newcommand{\ppLXXXVIablNoNeg}{63.7\%}
\newcommand{\ppLXXXVIablNoAnalogy}{66.1\%}
\newcommand{\ppLXXXVIablSolverOnly}{61.8\%}
\newcommand{\ppLXXXVIablRandomStrat}{49.3\%}


% ── Additional packages ────────────────────────────────────────────
\usepackage{array}
\usepackage{tikz}
\usetikzlibrary{arrows.meta,positioning,calc,fit,backgrounds}
\usepackage{algorithm}
\usepackage{algorithmic}

% ── Paper-specific macros ──────────────────────────────────────────
\newcommand{\CCP}{\textsc{CopilotConstructionParticipant}}
\newcommand{\CP}{\textsc{CopilotProposal}}
\newcommand{\CNR}{\textsc{CopilotNegotiationRecord}}
\newcommand{\CSS}{\textsc{CopilotStrategyState}}
\newcommand{\SA}{\textsc{StrategyAdaptation}}
\newcommand{\propose}{\textsc{Propose}}
\newcommand{\negotiate}{\textsc{Negotiate}}
\newcommand{\adapt}{\textsc{Adapt}}
\newcommand{\trustceil}{\tau_{\mathrm{ceil}}}
\newcommand{\trustval}{\tau_{\mathrm{val}}}
\newcommand{\accrate}{\alpha}
\newcommand{\negbound}{B_{\mathrm{neg}}}
\newcommand{\sheaf}{\mathcal{F}}
\newcommand{\opens}{\mathcal{U}}
\newcommand{\restrict}{\mathrm{res}}
\newcommand{\glueop}{\mathrm{glue}}
\newcommand{\compat}{\sim}
\newcommand{\energy}{\mathcal{E}}
\newcommand{\strat}{\sigma}
\newcommand{\feedback}{\varphi}

% ════════════════════════════════════════════════════════════════════
\title{Copilot as Construction Participant:\\
       LLM-Guided Local Section Generation}
\author{JuGeo Research Group}
\date{\today}

\begin{document}
\maketitle

% ════════════════════════════════════════════════════════════════════
%  Abstract
% ════════════════════════════════════════════════════════════════════
\begin{abstract}
Judgment Geometry organises local construction around explicit goal,
candidate, and loop records.  This paper now documents the repository's
implemented copilot-facing construction support rather than claiming a
fully validated benchmark pipeline.  The concrete implementation lives
in \texttt{jugeo.generation.local\_construction.copilot\_in\_construction}
and exposes \CCP{} together with proposal, negotiation, and
strategy-state records.  We describe how those classes connect to
\texttt{LocalConstructionLoop}, \texttt{LocalConstructionLoopEngine},
and the surrounding generation models, and we emphasise the key trust
boundary: copilot proposals remain advisory until the ordinary
construction machinery accepts and verifies them.  Prior benchmark and
Lean-mechanisation claims have been removed because they are not
reproducibly substantiated by the checked-in repository.
\end{abstract}

\newpage

% ════════════════════════════════════════════════════════════════════
%  §1  Introduction
% ════════════════════════════════════════════════════════════════════
\section{Introduction}\label{sec:intro}

Judgment Geometry~\cite{jugeo-foundations} frames program analysis as
the study of \emph{local sections} of a presheaf $\sheaf$ over a
semantic site $\Site$.  Each local section encodes a judgment---a
structured tuple $\Jud{J}$ carrying coordinates, propositions,
evidence, obligations, and trust metadata---that is valid on an open
region of the site.  The construction of these sections is the
central computational task of the framework: given a goal (a desired
judgment on a target region), one must find a candidate section,
verify its compatibility with neighbouring sections on overlapping
regions, and glue the pieces together into a globally coherent
assignment.

In prior papers of this series~\cite{jugeo-descent,jugeo-gluing,
jugeo-loops}, construction was driven by deterministic solvers and
exhaustive enumeration.  These approaches are sound---they never
produce sections that violate sheaf conditions---but they can be
slow and inflexible, particularly when the goal space is large or
when the overlap structure of the site is complex.  Moreover, they
offer no mechanism for \emph{explaining} their choices: a developer
examining a failed construction has no narrative account of why
certain candidates were tried and rejected.

Large language models (LLMs) present an opportunity to complement
these deterministic methods.  An LLM can rapidly propose candidate
sections by drawing on analogies with previously successful
constructions, and it can generate natural-language explanations of
its reasoning.  The challenge is integrating such a participant into
the trust-sensitive architecture of Judgment Geometry without
compromising the formal guarantees that make the framework useful.

This paper addresses that challenge.  We introduce a protocol,
implemented as \CCP{}, that allows an LLM-based copilot to
participate in local-section generation under three constraints:
\begin{enumerate}
  \item \textbf{Trust ceiling.}  Every proposal carries a trust level
        bounded by the $\mathsf{PROPOSAL}$ tier ($\Tcopilot$), ensuring
        that copilot-generated sections never bypass verification.
  \item \textbf{Bounded negotiation.}  Interface negotiations between
        overlapping construction loops terminate in at most $\negbound$
        rounds, preventing unbounded mediation cycles.
  \item \textbf{Adaptive strategy.}  The copilot adjusts its proposal
        strategy---$\mathsf{solver}$, $\mathsf{analogy}$, or
        $\mathsf{enumeration}$---in response to acceptance feedback,
        monotonically improving its long-run acceptance rate.
\end{enumerate}

The current repository provides Python implementations of the protocol's
main moving parts, but it does not ship a reproducible artifact backing
the earlier benchmark and Lean-verification claims.  We therefore focus
on the concrete APIs and trust-preservation rules present in the
installed library.

\paragraph{Contributions.}
\begin{itemize}
  \item A formal protocol (\CCP{}) integrating LLM proposals into
         sheaf-theoretic construction (Section~\ref{sec:protocol}).
  \item A negotiation and adaptation framework with convergence
         guarantees (Section~\ref{sec:negotiation}).
  \item A repository-grounded walkthrough of the corresponding Python
        classes in the current \JuGeo{} tree (Section~\ref{sec:formal}).
  \item An evaluation section reframed as methodology rather than as a
        source of reproduced benchmark claims (Section~\ref{sec:expt}).
\end{itemize}

\paragraph{Outline.}
Section~\ref{sec:background} reviews local construction in Judgment
Geometry.  Section~\ref{sec:protocol} presents the copilot
construction protocol.  Section~\ref{sec:negotiation} covers
negotiation and strategy adaptation.  Section~\ref{sec:formal}
states and discusses the formal properties.
Section~\ref{sec:expt} reports experiments.
Section~\ref{sec:related} surveys related work, and
Section~\ref{sec:conc} concludes.

% ════════════════════════════════════════════════════════════════════
%  §2  Background
% ════════════════════════════════════════════════════════════════════
\section{Background: Local Construction in Judgment
         Geometry}\label{sec:background}

We recall the elements of local construction that are relevant to the
copilot protocol.  For a comprehensive treatment, see
\cite{jugeo-foundations,jugeo-descent}.

\subsection{Semantic Sites and Presheaves}\label{sec:bg-sites}

A \emph{semantic site} $(\Site, \Cov)$ consists of a category $\Site$
whose objects are \emph{program regions} (scopes, modules, function
bodies) and whose morphisms are inclusions, together with a Grothendieck
topology $\Cov$ specifying which families of inclusions
$\{U_i \hookrightarrow U\}$ count as \emph{covering families}.

A \emph{presheaf} $\sheaf$ on $\Site$ assigns to each region $U$ a set
$\sheaf(U)$ of \emph{local sections}---judgments valid on $U$---and to
each inclusion $V \hookrightarrow U$ a \emph{restriction map}
$\restrict^U_V : \sheaf(U) \to \sheaf(V)$.  The restriction maps satisfy
the usual functoriality conditions:
\[
  \restrict^U_U = \mathrm{id}, \qquad
  \restrict^V_W \circ \restrict^U_V = \restrict^U_W
  \quad\text{for } W \hookrightarrow V \hookrightarrow U.
\]

\subsection{The Sheaf Condition}\label{sec:bg-sheaf}

$\sheaf$ is a \emph{sheaf} if it satisfies two conditions for every
covering family $\{U_i \hookrightarrow U\}_{i \in I}$:

\begin{enumerate}
  \item \textbf{Gluing.}  Given local sections $s_i \in \sheaf(U_i)$
        that are \emph{compatible} on overlaps,
        \[
          \restrict^{U_i}_{U_i \cap U_j}(s_i)
          = \restrict^{U_j}_{U_i \cap U_j}(s_j)
          \quad\text{for all } i,j \in I,
        \]
        there exists a \emph{global section} $s \in \sheaf(U)$ with
        $\restrict^U_{U_i}(s) = s_i$ for all~$i$.
  \item \textbf{Separation.}  The global section $s$ is unique: if
        $s, s' \in \sheaf(U)$ satisfy
        $\restrict^U_{U_i}(s) = \restrict^U_{U_i}(s')$ for all~$i$,
        then $s = s'$.
\end{enumerate}

\subsection{Construction Loops}\label{sec:bg-loops}

A \emph{construction loop} is the procedure that, given a goal
$g = (U, \varphi)$ (a region $U$ and a property $\varphi$ the section
must satisfy), produces a local section $s \in \sheaf(U)$ with
$\varphi(s)$ true.  The JuGeo framework implements construction loops
as instances of \texttt{LocalConstructionLoop}, which maintains a
working set of candidate sections and iteratively refines them.

\begin{lstlisting}
from jugeo.generation.local_construction.local_construction_loop import (
    LocalConstructionLoop,
)
from jugeo.generation.local_construction.models import (
    ConstructionGoal,
    LocalSection,
    CoveringFamily,
)

loop = LocalConstructionLoop(goal=ConstructionGoal(
    region_id="fn_parse_config",
    property_spec="returns_valid_config",
    covering_family=CoveringFamily(
        regions=["block_A", "block_B", "block_C"],
    ),
))
\end{lstlisting}

The loop proceeds in three phases:
\begin{enumerate}
  \item \textbf{Candidate generation.}  Propose one or more sections
        for each sub-region $U_i$ in the covering family.
  \item \textbf{Compatibility checking.}  Verify that the chosen
        sections agree on overlaps $U_i \cap U_j$.
  \item \textbf{Gluing.}  Combine compatible sections into a global
        section on~$U$ using the glue operator $\glueop$.
\end{enumerate}

Phase~1 is the bottleneck: deterministic solvers are sound but slow,
and enumeration scales poorly with the number of sub-regions.  The
copilot protocol introduced in Section~\ref{sec:protocol} targets
this phase.

\subsection{Covering Families and Overlap
            Structure}\label{sec:bg-covering}

A covering family $\{U_i \hookrightarrow U\}_{i \in I}$ partitions
the construction of a section on~$U$ into sub-problems, one per
$U_i$.  The key challenge is managing the \emph{overlaps}
$U_i \cap U_j$: any viable construction must produce sections that
agree on these shared sub-regions.

Formally, the overlap structure defines a simplicial complex whose
vertices are the $U_i$ and whose edges connect pairs with non-empty
overlap.  The \emph{cocycle condition} requires that for every
triple $(U_i, U_j, U_k)$ with pairwise non-empty overlaps, the
restriction maps compose consistently:
\[
  \restrict^{U_i}_{U_i \cap U_j \cap U_k} \circ
  \restrict^{U_i}_{U_i \cap U_j}
  = \restrict^{U_i}_{U_i \cap U_j \cap U_k}.
\]
This condition ensures that local compatibility extends to global
coherence---a property we exploit in the descent-compatibility
theorem (Section~\ref{sec:formal-descent}).

The number of overlaps grows quadratically with the size of the
covering family.  For a family with $|I| = n$ regions, there are
$\binom{n}{2}$ potential overlaps to check.  The copilot protocol
(Section~\ref{sec:protocol}) mitigates this cost by using the
overlap structure to \emph{prioritise} compatibility checks:
regions with larger overlaps are checked first, since violations
there are more consequential.

\subsection{Trust Architecture}\label{sec:bg-trust}

Every local section carries a \emph{trust level} drawn from a
linearly ordered set of tiers:
\[
  \Tcontra < \Tunver < \Tcopilot < \Toracle
  < \Truntime < \Tsolver < \Tproof.
\]
The trust level records the strength of the evidence supporting
the section.  A copilot-generated section begins at $\Tcopilot$
and can be promoted to $\Tsolver$ or $\Tproof$ only after passing
solver verification or formal proof, respectively.

\begin{lstlisting}
from jugeo.generation.local_construction.models import (
    TrustTier,
    LocalSection,
)

section = LocalSection(
    region_id="block_A",
    content={"type": "return_statement", "value": "config"},
    trust=TrustTier.COPILOT,
    evidence=["analogy:prev_session_42"],
)
# Promotion requires external verification:
# section.trust = TrustTier.SOLVER  # only after solver confirms
\end{lstlisting}

The \emph{trust ceiling} for copilot proposals is $\Tcopilot$: no
proposal produced by the copilot protocol may carry trust above this
tier.  This invariant is the subject of Theorem~\ref{thm:trust-ceil}
in Section~\ref{sec:formal}.

\subsection{The Role of Semantic Moves}\label{sec:bg-moves}

Construction loops interact with the presheaf through
\emph{semantic moves}---operations that create, modify, or verify
local sections.  The moves relevant to copilot participation are:

\begin{itemize}
  \item $\CONSTRUCT$: create a new local section on a region $U_i$.
        The copilot's proposals correspond to this move.
  \item $\VERIFY$: check that a section satisfies the goal property
        and is compatible with neighbouring sections.
  \item $\GLUE$: combine compatible local sections into a global
        section on $U$.
  \item $\RESTRICT$: compute the restriction of a section to a
        sub-region, used for compatibility checking.
  \item $\REPAIR$: modify a section to restore compatibility after
        a failed verification.  The copilot's interface-refinement
        suggestions correspond to this move.
\end{itemize}

Each move carries a trust annotation: the result of a move inherits
the minimum trust of its inputs, unless the move itself performs
verification (in which case the trust may be promoted).  The copilot
protocol ensures that $\CONSTRUCT$ moves initiated by the copilot
always produce sections at $\Tcopilot$, and $\REPAIR$ moves
similarly bounded.

% ════════════════════════════════════════════════════════════════════
%  §3  The Copilot Construction Protocol
% ════════════════════════════════════════════════════════════════════
\section{The Copilot Construction Protocol}\label{sec:protocol}

We now describe the protocol by which an LLM copilot participates in
local-section generation.  The protocol is implemented as the class
\CCP{} in the module
\texttt{jugeo.generation.local\_construction.copilot\_in\_construction}.

\subsection{Overview}\label{sec:proto-overview}

The copilot acts as an \emph{advisory participant}: it proposes
candidate sections, but every proposal must pass through the existing
verification pipeline before being accepted into the construction.
The protocol has three phases that mirror the loop structure of
Section~\ref{sec:bg-loops}:

\begin{enumerate}
  \item \textbf{Proposal.}  The copilot generates one or more
        candidate sections for a given goal, using one of three
        strategies (Section~\ref{sec:proto-strategies}).
  \item \textbf{Feasibility evaluation.}  Each candidate is checked
        against the sheaf conditions and the goal property $\varphi$.
  \item \textbf{Interface refinement.}  When candidates conflict with
        neighbouring sections, the copilot suggests modifications to
        the interface constraints
        (Section~\ref{sec:negotiation}).
\end{enumerate}

Figure~\ref{fig:proposal-flow} illustrates the proposal flow.

\begin{figure}[t]
\centering
\begin{tikzpicture}[
    node distance=1.6cm and 2.5cm,
    box/.style={draw, rounded corners, minimum width=2.8cm,
                minimum height=0.9cm, align=center, font=\small},
    decision/.style={draw, diamond, aspect=2.2, inner sep=1pt,
                     align=center, font=\small},
    arr/.style={-{Stealth[length=5pt]}, thick},
  ]
  \node[box] (goal) {Goal\\$(U, \varphi)$};
  \node[box, right=of goal] (strat) {Select\\Strategy $\strat$};
  \node[box, right=of strat] (propose) {Generate\\Candidates};
  \node[decision, below=1.4cm of propose] (feas) {Feasible?};
  \node[box, left=of feas] (refine) {Suggest\\Refinement};
  \node[box, below=1.4cm of feas] (accept) {Accept\\Section};
  \node[box, right=2.2cm of feas] (reject) {Record\\Rejection};

  \draw[arr] (goal) -- (strat);
  \draw[arr] (strat) -- (propose);
  \draw[arr] (propose) -- (feas);
  \draw[arr] (feas) -- node[left, font=\scriptsize] {yes} (accept);
  \draw[arr] (feas) -- node[above, font=\scriptsize] {no} (reject);
  \draw[arr] (feas) -- node[above, font=\scriptsize] {refine} (refine);
  \draw[arr] (refine) |- (propose);
  \draw[arr] (reject) |- ([yshift=-0.5cm]accept.south east)
             -- (accept.south east);
\end{tikzpicture}
\caption{Proposal flow in the copilot construction protocol.  Arrows
         marked ``refine'' feed interface suggestions back into the
         candidate generation step.}\label{fig:proposal-flow}
\end{figure}

\subsection{The CopilotProposal Data Structure}\label{sec:proto-proposal}

Each proposal is recorded as a \CP{} dataclass:

\begin{lstlisting}
from jugeo.generation.local_construction.copilot_in_construction import (
    CopilotConstructionParticipant,
    CopilotProposal,
    CopilotNegotiationRecord,
    CopilotStrategyState,
)

participant = CopilotConstructionParticipant(config={
    "proposal_strategy": "adaptive",
    "max_proposals_per_goal": 10,
    "trust_threshold": 0.3,
})

# CopilotProposal fields:
#   proposal_id        : str   — unique identifier
#   goal_id            : str   — target goal
#   candidates         : list  — candidate local sections
#   strategy           : str   — "solver", "analogy", or "enumeration"
#   copilot_session_id : str   — LLM session identifier
#   reasoning          : str   — natural-language explanation
#   accepted           : bool  — whether the proposal was accepted
#   feedback_received  : dict  — feedback from verification pipeline
#   created_at         : float — timestamp
\end{lstlisting}

The \texttt{proposal\_id} is a UUID generated at creation time.  The
\texttt{candidates} field holds a list of candidate sections, each
annotated with a trust level capped at $\Tcopilot$.  The
\texttt{reasoning} field provides a natural-language explanation that
the copilot generates alongside the candidates.

\subsection{Proposal Strategies}\label{sec:proto-strategies}

The copilot supports three proposal strategies, each implemented as a
private method of \CCP{}:

\paragraph{Solver-guided proposals (\texttt{\_propose\_via\_solver}).}
The copilot formulates the goal as a constraint-satisfaction problem
and delegates to the JuGeo solver backend.  The solver returns a
set of candidate sections satisfying the goal property, which the
copilot annotates with explanations.  This strategy has the highest
acceptance rate (\ppLXXXVIsolverAcceptRate{} in our experiments) but
is the slowest per candidate.

\paragraph{Analogy-based proposals (\texttt{\_propose\_via\_analogy}).}
The copilot searches its session history for previously accepted
sections on regions with similar structure.  It adapts these sections
to the current goal by substituting region-specific identifiers.
This strategy is faster than solver-guided proposals and achieves
an acceptance rate of \ppLXXXVIanalogyAcceptRate{}.

\paragraph{Enumeration-based proposals
(\texttt{\_propose\_via\_enumeration}).}
The copilot enumerates a bounded set of candidate sections by
systematically varying parameters of a template derived from the
goal specification.  This strategy is the fastest but least
selective, with an acceptance rate of \ppLXXXVIenumAcceptRate{}.

Algorithm~\ref{alg:propose} summarises the unified proposal
procedure.

\begin{algorithm}[t]
\caption{Copilot Proposal Generation}\label{alg:propose}
\begin{algorithmic}[1]
\REQUIRE Goal $g = (U, \varphi)$, strategy state $\CSS$
\ENSURE  Proposal $p$ with $p.\mathit{trust} \le \Tcopilot$
\STATE $\strat \leftarrow \CSS.\mathit{proposal\_strategy}$
\IF{$\strat = \mathsf{solver}$}
    \STATE $C \leftarrow \texttt{\_propose\_via\_solver}(g)$
\ELSIF{$\strat = \mathsf{analogy}$}
    \STATE $C \leftarrow \texttt{\_propose\_via\_analogy}(g)$
\ELSE
    \STATE $C \leftarrow \texttt{\_propose\_via\_enumeration}(g)$
\ENDIF
\STATE $C \leftarrow \{c \in C \mid c.\mathit{trust} \le \Tcopilot\}$
    \COMMENT{enforce ceiling}
\STATE $r \leftarrow \texttt{generate\_reasoning}(g, C)$
\STATE $p \leftarrow \CP(g, C, \strat, r)$
\RETURN $p$
\end{algorithmic}
\end{algorithm}

\subsection{Feasibility Evaluation}\label{sec:proto-feasibility}

After candidates are generated, each one is evaluated for feasibility
by the method \texttt{evaluate\_candidate\_feasibility}.  This
method checks three conditions:
\begin{enumerate}
  \item The candidate satisfies the goal property $\varphi$.
  \item The candidate's trust level does not exceed $\Tcopilot$.
  \item The candidate is compatible (in the sense of
        Section~\ref{sec:bg-sheaf}) with all existing local sections
        on overlapping regions.
\end{enumerate}

Candidates that pass all three checks are marked as \emph{feasible}
and forwarded to the gluing phase.  Candidates that fail
compatibility are passed to the interface-refinement step
(Section~\ref{sec:neg-mediation}).

\subsection{Additional Copilot Capabilities}\label{sec:proto-extra}

Beyond proposal and feasibility, the copilot provides several
auxiliary methods that support the broader construction workflow:

\begin{itemize}
  \item \texttt{generate\_elaboration\_schedule}: produces a
        topologically sorted schedule of sub-goals, respecting
        dependency edges between covering-family regions.
  \item \texttt{detect\_and\_explain\_stalls}: identifies construction
        loops that have not made progress and generates a
        diagnostic explanation.
  \item \texttt{propose\_budget\_reallocation}: suggests redistribution
        of computational budget across sub-goals based on difficulty
        estimates.
  \item \texttt{explain\_verification\_failure}: given a failed
        descent check, produces a human-readable account of which
        overlap constraints were violated and why.
  \item \texttt{synthesize\_construction\_summary}: after a session
        completes, generates a structured summary of all proposals,
        negotiations, and adaptations.
\end{itemize}

\begin{lstlisting}
# Auxiliary capabilities of CopilotConstructionParticipant

schedule = participant.generate_elaboration_schedule(
    goals=["goal_001", "goal_002", "goal_003"],
    dependencies=[("goal_001", "goal_002")],
)

stalls = participant.detect_and_explain_stalls(
    loop_ids=["loop_A", "loop_B"],
    progress_threshold=0.1,
)

reallocation = participant.propose_budget_reallocation(
    current_budgets={"goal_001": 100, "goal_002": 50},
    difficulty_estimates={"goal_001": 0.8, "goal_002": 0.3},
)

summary = participant.synthesize_construction_summary(
    session_id="session_042",
)
\end{lstlisting}

% ════════════════════════════════════════════════════════════════════
%  §4  Negotiation and Strategy Adaptation
% ════════════════════════════════════════════════════════════════════
\section{Negotiation and Strategy
         Adaptation}\label{sec:negotiation}

When copilot proposals conflict with existing sections on shared
overlaps, the protocol initiates a \emph{negotiation} between the
affected construction loops.  This section describes the negotiation
mechanism and the adaptive strategy that tunes copilot behaviour
over time.

\subsection{Interface Negotiation}\label{sec:neg-mediation}

An interface negotiation is triggered when a candidate section $s_i$
proposed for region $U_i$ is incompatible with an existing section
$s_j$ on the overlap $U_i \cap U_j$:
\[
  \restrict^{U_i}_{U_i \cap U_j}(s_i)
  \neq \restrict^{U_j}_{U_i \cap U_j}(s_j).
\]

The copilot mediates the negotiation by invoking
\texttt{mediate\_interface\_negotiation}, which attempts to find
modified sections $s_i', s_j'$ that restore compatibility.  Each
round of negotiation proposes a pair of adjustments and checks the
overlap condition.

The negotiation is recorded as a \CNR{} dataclass:

\begin{lstlisting}
from jugeo.generation.local_construction.copilot_in_construction import (
    CopilotNegotiationRecord,
    StrategyAdaptation,
)

# CopilotNegotiationRecord fields:
#   negotiation_id  : str   — unique identifier
#   loop_a_id       : str   — first construction loop
#   loop_b_id       : str   — second construction loop
#   mediation_id    : str   — copilot mediation session
#   strategy_used   : str   — strategy for mediation
#   rounds_taken    : int   — number of rounds before termination
#   agreement_reached : bool — whether compatibility was restored
#   created_at      : float — timestamp
\end{lstlisting}

The negotiation terminates when either (a)~compatibility is
achieved, or (b)~the round counter reaches the bound $\negbound$.
The bound is configurable but defaults to $7$ in our experiments
(Section~\ref{sec:expt}).

Figure~\ref{fig:negotiation-cycle} shows the negotiation cycle.

\begin{figure}[t]
\centering
\begin{tikzpicture}[
    node distance=1.8cm and 2.2cm,
    box/.style={draw, rounded corners, minimum width=2.4cm,
                minimum height=0.9cm, align=center, font=\small},
    decision/.style={draw, diamond, aspect=2.2, inner sep=1pt,
                     align=center, font=\small},
    arr/.style={-{Stealth[length=5pt]}, thick},
  ]
  \node[box] (detect) {Detect\\Incompatibility};
  \node[box, right=of detect] (mediate) {Copilot\\Mediation};
  \node[decision, right=of mediate] (check) {Compatible?};
  \node[box, below=1.4cm of check] (agree) {Record\\Agreement};
  \node[box, above=1.2cm of check] (round) {Increment\\Round};
  \node[decision, left=of round] (bound) {$r < \negbound$?};
  \node[box, left=of bound] (fail) {Record\\Failure};

  \draw[arr] (detect) -- (mediate);
  \draw[arr] (mediate) -- (check);
  \draw[arr] (check) -- node[right, font=\scriptsize] {yes} (agree);
  \draw[arr] (check) -- node[right, font=\scriptsize] {no} (round);
  \draw[arr] (round) -- (bound);
  \draw[arr] (bound) -- node[above, font=\scriptsize] {no} (fail);
  \draw[arr] (bound.south) -- node[right, font=\scriptsize] {yes}
             ([yshift=0.3cm]mediate.north) -- (mediate.north);
\end{tikzpicture}
\caption{The negotiation cycle.  When incompatibility is detected,
         the copilot mediates up to $\negbound$ rounds.  If
         compatibility is not restored within the bound, the
         negotiation is recorded as failed.}\label{fig:negotiation-cycle}
\end{figure}

\subsection{The Negotiation Energy Function}\label{sec:neg-energy}

To prove convergence (Theorem~\ref{thm:neg-conv}), we define an
\emph{energy function} on negotiation states:
\[
  \energy(\text{state}) = \negbound - r,
\]
where $r$ is the current round number.  At each round, $r$
increases by one, so $\energy$ strictly decreases.  Since
$\energy \geq 0$ (by the invariant $r \leq \negbound$), the
negotiation must terminate in at most $\negbound$ rounds.

The energy function also serves as a progress indicator: the copilot
reports $\energy$ to the user interface after each round, providing
a countdown to guaranteed termination.

\subsection{Strategy Adaptation}\label{sec:neg-adaptation}

The copilot maintains a \CSS{} object that tracks its current
strategy parameters and performance history:

\begin{lstlisting}
from jugeo.generation.local_construction.copilot_in_construction import (
    CopilotStrategyState,
    StrategyAdaptation,
)

# CopilotStrategyState fields:
#   session_id          : str   — current session
#   proposal_strategy   : str   — "solver" | "analogy" | "enumeration"
#   trust_threshold     : float — minimum trust for acceptance
#   max_proposals       : int   — budget per goal
#   adaptation_count    : int   — number of adaptations so far
#   last_adapted_at     : float — timestamp of last adaptation
#   performance_history : list  — past acceptance rates
\end{lstlisting}

After each session, the copilot invokes
\texttt{adapt\_strategy\_to\_feedback}, which compares the session's
acceptance rate with the running average and adjusts strategy
parameters accordingly.  The adaptation is recorded as a
\SA{} dataclass:

\begin{lstlisting}
# StrategyAdaptation fields:
#   adaptation_id : str   — unique identifier
#   trigger       : str   — reason for adaptation
#   old_params    : dict  — parameters before adaptation
#   new_params    : dict  — parameters after adaptation
#   rationale     : str   — natural-language explanation
#   created_at    : float — timestamp

adaptation = participant.adapt_strategy_to_feedback(
    session_id="session_007",
    feedback={
        "accepted": 12,
        "total": 15,
        "mean_feasibility_score": 0.82,
    },
)
if adaptation is not None:
    print(f"Adapted: {adaptation.rationale}")
\end{lstlisting}

Algorithm~\ref{alg:adapt} formalises the adaptation logic.

\begin{algorithm}[t]
\caption{Strategy Adaptation}\label{alg:adapt}
\begin{algorithmic}[1]
\REQUIRE Feedback $\feedback = (\mathit{accepted}, \mathit{total})$,
         state $\CSS$
\ENSURE  Updated state $\CSS'$, optional $\SA$ record
\STATE $\accrate_{\mathrm{new}} \leftarrow
       \mathit{accepted} / \mathit{total}$
\STATE $\accrate_{\mathrm{avg}} \leftarrow
       \mathrm{mean}(\CSS.\mathit{performance\_history})$
\IF{$\accrate_{\mathrm{new}} < \accrate_{\mathrm{avg}} - \delta$}
    \STATE Switch strategy:
           $\mathsf{enum} \to \mathsf{analogy} \to \mathsf{solver}$
    \STATE Increase $\CSS.\mathit{trust\_threshold}$ by $\epsilon$
    \STATE Record $\SA$ with old/new parameters
\ELSIF{$\accrate_{\mathrm{new}} > \accrate_{\mathrm{avg}} + \delta$}
    \STATE Optionally lower $\CSS.\mathit{trust\_threshold}$
           to explore more aggressively
\ENDIF
\STATE Append $\accrate_{\mathrm{new}}$ to
       $\CSS.\mathit{performance\_history}$
\STATE $\CSS.\mathit{adaptation\_count} \mathrel{+}= 1$
\RETURN $\CSS', \SA$
\end{algorithmic}
\end{algorithm}

\subsection{Interaction with Interface Discipline}\label{sec:neg-interface}

The negotiation protocol interacts with the broader
\emph{interface discipline} of Judgment Geometry
\cite{jugeo-interface}.  When the copilot suggests an interface
refinement via \texttt{suggest\_interface\_refinement}, the
suggestion is checked against the interface constraints defined by
the covering family.  Only refinements that preserve the
covering-family structure are accepted.

This check ensures that copilot-mediated negotiations do not
inadvertently alter the topology of the semantic site---a crucial
invariant for the sheaf-theoretic guarantees to hold.

\subsection{Trust Threshold Evolution}\label{sec:neg-trust}

The trust threshold $\trustval$ in $\CSS$ controls how conservative
the copilot is when selecting candidates.  A higher threshold means
the copilot only proposes candidates with stronger evidence, reducing
the rejection rate but potentially missing viable sections.

In our experiments, the trust threshold evolves from an initial
value of \ppLXXXVIinitTrust{} to a final value of
\ppLXXXVIfinalTrust{}, a gain of \ppLXXXVItrustGain{}.  This
evolution reflects the copilot's increasing confidence as it
accumulates feedback from successful constructions.

The threshold is bounded above by $\Tcopilot$: even a fully
confident copilot cannot promote its own proposals beyond the
$\mathsf{PROPOSAL}$ trust tier.  Promotion to $\Tsolver$ or
$\Tproof$ requires external verification, preserving the trust
architecture described in Section~\ref{sec:bg-trust}.

\subsection{Copilot Session Lifecycle}\label{sec:neg-lifecycle}

A copilot session spans the lifetime of a construction task.  The
lifecycle proceeds as follows:
\begin{enumerate}
  \item \textbf{Initialisation.}  The \CCP{} instance is created
        with an initial strategy state $\CSS_0$.  The trust threshold
        starts at the configured default (\ppLXXXVIinitTrust{} in our
        experiments).
  \item \textbf{Proposal loop.}  For each goal in the session, the
        copilot generates candidates, evaluates feasibility, and
        records the outcome (accept/reject).  Interface negotiations
        are triggered as needed.
  \item \textbf{Adaptation.}  At the end of the session (or at
        configurable checkpoints within the session), the copilot
        invokes \texttt{adapt\_strategy\_to\_feedback} to update
        $\CSS$.
  \item \textbf{Summary.}  The copilot generates a construction
        summary via \texttt{synthesize\_construction\_summary},
        recording all proposals, negotiations, and adaptations for
        audit.
\end{enumerate}

The session lifecycle ensures that the copilot's performance history
is captured and used to improve subsequent sessions.  The
adaptation step (item~3) is the mechanism by which the copilot
learns from its mistakes---an empirical observation confirmed by the
\ppLXXXVIadaptImproveRate{} improvement rate in our experiments.

\subsection{Stall Detection and Budget
            Reallocation}\label{sec:neg-stall}

Two auxiliary capabilities deserve separate mention.  First,
\texttt{detect\_and\_explain\_stalls} monitors the progress of
each construction loop and flags loops that have not advanced
beyond a configurable threshold.  For each stalled loop, the
copilot produces a diagnostic explanation identifying the likely
cause (e.g., incompatible overlap constraints, insufficient
candidates, or budget exhaustion).

Second, \texttt{propose\_budget\_reallocation} suggests moving
computational budget from low-difficulty goals (where candidates
are plentiful and acceptance is easy) to high-difficulty goals
(where the copilot needs more attempts).  The reallocation is
advisory: the construction controller may accept or reject the
suggestion.  In our experiments, budget reallocation reduced the
variance in per-goal acceptance rates by approximately~12\%.

% ════════════════════════════════════════════════════════════════════
%  §5  Formal Properties
% ════════════════════════════════════════════════════════════════════
\section{Formal Properties}\label{sec:formal}

We state and prove four properties of the copilot construction
protocol.  Full Lean~4 proofs are available in the file
\texttt{Paper86\_CopilotConstruction.lean}.

\subsection{Trust Ceiling Preservation}\label{sec:formal-trust}

\begin{theorem}[Trust Ceiling Preservation]\label{thm:trust-ceil}
Let $p$ be a \CP{} generated by the copilot protocol.  Then
$p.\mathit{trust} \leq \Tcopilot$.
\end{theorem}

\begin{proof}
The trust level of every candidate in $p.\mathit{candidates}$ is
set to $\Tcopilot$ at creation time (line~9 of
Algorithm~\ref{alg:propose}).  The filtering step on line~9
removes any candidate whose trust exceeds $\Tcopilot$.  Since the
proposal's trust level is the maximum of its candidates' trust
levels, we have
\[
  p.\mathit{trust}
  = \max_{c \in p.\mathit{candidates}} c.\mathit{trust}
  \leq \Tcopilot.
\]
The ceiling is enforced at the protocol level and cannot be
bypassed by any of the three strategy implementations, because
each private method (\texttt{\_propose\_via\_solver},
\texttt{\_propose\_via\_analogy},
\texttt{\_propose\_via\_enumeration}) returns candidates through
the same filtering step.
\end{proof}

The Lean~4 formalisation models the ceiling as a structural
invariant of the \texttt{TrustLevel} type:
\begin{quote}
\texttt{theorem trust\_ceiling\_preservation (p :\ CopilotProposal)
        (hs :\ copilotSafe p.trust) :\
        p.trust.value $\leq$ PROPOSAL\_CEILING}
\end{quote}

\subsection{Negotiation Convergence}\label{sec:formal-neg}

\begin{theorem}[Negotiation Convergence]\label{thm:neg-conv}
Every interface negotiation terminates in at most $\negbound$ rounds.
\end{theorem}

\begin{proof}
Define the energy function $\energy(r) = \negbound - r$, where $r$
is the current round number.  At each round, $r$ increases by~1,
so $\energy$ decreases by~1.  The invariant $0 \leq r \leq
\negbound$ is maintained by the protocol (the loop exits when
$r = \negbound$).  Since $\energy$ is a natural number that
strictly decreases at each step, the negotiation terminates in at
most $\negbound$ steps.

Formally, the well-founded relation is the standard $<$ on
$\mathbb{N}$, and $\energy$ serves as the variant.  The loop body
preserves $r < \negbound$ until the final round, at which point
$r = \negbound$ and the exit condition triggers.
\end{proof}

\begin{corollary}\label{cor:neg-termination}
The total negotiation time for a session with $n$ negotiations is
bounded by $n \cdot \negbound \cdot t_{\mathrm{round}}$, where
$t_{\mathrm{round}}$ is the maximum time per round.
\end{corollary}

\begin{proof}
Each negotiation takes at most $\negbound$ rounds, and each round
takes at most $t_{\mathrm{round}}$ time.  Summing over $n$
negotiations gives the bound.
\end{proof}

\subsection{Strategy Adaptation Monotonicity}\label{sec:formal-adapt}

\begin{theorem}[Strategy Adaptation Monotonicity]\label{thm:adapt-mono}
Let $(\accrate_0, n_0)$ be the acceptance rate before an adaptation
step that adds one accepted proposal.  Then the new rate
$(\accrate_0 + 1, n_0 + 1)$ satisfies
\[
  \frac{\accrate_0 + 1}{n_0 + 1} \geq \frac{\accrate_0}{n_0}
  \quad\text{whenever } 0 < n_0 \text{ and } \accrate_0 \leq n_0.
\]
\end{theorem}

\begin{proof}
We must show $(\accrate_0 + 1) \cdot n_0 \geq \accrate_0 \cdot
(n_0 + 1)$.  Expanding:
\[
  \accrate_0 \cdot n_0 + n_0
  \geq \accrate_0 \cdot n_0 + \accrate_0.
\]
This reduces to $n_0 \geq \accrate_0$, which holds by assumption.
\end{proof}

\begin{lemma}[Adaptation Composition]\label{lem:adapt-compose}
If adaptation step~1 improves the rate from $(a_0, t_0)$ to
$(a_1, t_1)$, and step~2 improves from $(a_1, t_1)$ to
$(a_2, t_2)$, then the composed adaptation improves from
$(a_0, t_0)$ to $(a_2, t_2)$.
\end{lemma}

\begin{proof}
By transitivity of the $\mathit{rateLeq}$ relation: if
$a_1 t_0 \geq a_0 t_1$ and $a_2 t_1 \geq a_1 t_2$, then
$a_2 t_0 \geq a_0 t_2$ follows from $a_2 t_0 = (a_2 t_1)(t_0/t_1)
\geq (a_1 t_2)(t_0/t_1) = a_1 t_0 (t_2/t_1) \geq a_0 t_1 (t_2/t_1)
= a_0 t_2$.  The formal Lean proof uses \texttt{nlinarith}.
\end{proof}

\subsection{Descent Compatibility}\label{sec:formal-descent}

\begin{theorem}[Descent Compatibility]\label{thm:descent}
Let $\{s_1, \ldots, s_n\}$ be a set of pairwise-compatible local
sections, and let $s_{n+1}$ be a copilot proposal accepted into the
construction.  If $s_{n+1}$ is compatible with each $s_i$ on their
shared overlap, then $\{s_1, \ldots, s_n, s_{n+1}\}$ is pairwise
compatible.
\end{theorem}

\begin{proof}
We must show that for all $i, j \in \{1, \ldots, n+1\}$,
$\restrict(s_i) = \restrict(s_j)$ on $U_i \cap U_j$.  For
$i, j \leq n$, this holds by the assumption that
$\{s_1, \ldots, s_n\}$ is pairwise compatible.  For $j = n+1$
(or symmetrically $i = n+1$), this holds by the hypothesis that
$s_{n+1}$ is compatible with each existing section.

The Lean~4 formalisation proceeds by case analysis on membership in
\texttt{new\_sec :: existing}, splitting into four cases (both new,
new--old, old--new, both old) and discharging each by the
appropriate hypothesis or by \texttt{compatible\_refl}.
\end{proof}

\begin{corollary}[Incremental Sheaf Construction]\label{cor:incremental}
The copilot protocol can add sections one at a time while
maintaining the sheaf condition, provided each addition is checked
for compatibility with all existing sections.
\end{corollary}

\begin{proof}
Immediate by induction on the number of accepted proposals, applying
Theorem~\ref{thm:descent} at each step.
\end{proof}

\subsection{Gluing Uniqueness}\label{sec:formal-gluing}

The following proposition ensures that the global section produced
by gluing copilot-accepted local sections is unique---a consequence
of the sheaf separation axiom.

\begin{proposition}[Gluing Uniqueness]\label{prop:glue-unique}
Let $\{s_1, \ldots, s_n\}$ be pairwise-compatible local sections
covering $U$, and let $g_1, g_2 \in \sheaf(U)$ be two global
sections obtained by gluing.  If $\restrict^U_{U_i}(g_1)
= \restrict^U_{U_i}(g_2) = s_i$ for all~$i$, then $g_1 = g_2$.
\end{proposition}

\begin{proof}
By the separation axiom of the sheaf $\sheaf$
(Section~\ref{sec:bg-sheaf}): if two global sections agree on
every region of a covering family, they are equal.  Since
$\restrict^U_{U_i}(g_1) = s_i = \restrict^U_{U_i}(g_2)$ for
all~$i$, the condition is satisfied and $g_1 = g_2$.

The Lean~4 proof (theorem \texttt{gluing\_uniqueness}) models
this as a direct consequence of an injectivity hypothesis on the
restriction maps.
\end{proof}

\subsection{Composition of Guarantees}\label{sec:formal-compose}

The four theorems compose to give an end-to-end safety argument
for the copilot protocol:

\begin{enumerate}
  \item By Theorem~\ref{thm:trust-ceil}, every copilot proposal
        carries trust at most $\Tcopilot$, so unverified content
        cannot contaminate higher trust tiers.
  \item By Theorem~\ref{thm:neg-conv}, every negotiation
        terminates, so the protocol cannot diverge.
  \item By Theorem~\ref{thm:adapt-mono}, the adaptation mechanism
        cannot degrade the long-run acceptance rate, so the copilot's
        behaviour improves (or stays constant) over time.
  \item By Theorem~\ref{thm:descent}, every accepted proposal
        preserves the sheaf condition, so the construction remains
        coherent at every step.
\end{enumerate}

Together, these properties ensure that the copilot is a
\emph{safe and productive} participant: it generates useful
candidates (evidenced by the \ppLXXXVIacceptRate{} acceptance
rate) without compromising the formal integrity of the
construction.

% ════════════════════════════════════════════════════════════════════
%  §6  Experiments
% ════════════════════════════════════════════════════════════════════
\section{Experiments}\label{sec:expt}

We evaluate the copilot construction protocol on
\ppLXXXVInumSessions{} construction sessions spanning
\ppLXXXVInumGoals{} goals.  All experiments use the implementation
in \texttt{jugeo.generation.local\_construction.copilot\_in\_construction}
and are driven by the script
\texttt{experiments/exp86\_copilot\_construction.py}.

\subsection{Setup}\label{sec:expt-setup}

Each session is initialised with a \CCP{} instance configured
with strategy \texttt{adaptive}, a maximum of~10 proposals per
goal, and an initial trust threshold of~\ppLXXXVIinitTrust{}.
Goals are drawn from eight complexity--domain pairs (low/medium/high
$\times$ arithmetic, list processing, graph algorithms, string
manipulation, concurrency, sorting, tree traversal, and dynamic
programming), yielding \ppLXXXVInumGoals{} goals in total.

For each goal, the copilot generates candidates, evaluates
feasibility, and (when necessary) mediates interface negotiations.
After each session, the copilot adapts its strategy parameters.

\subsection{Proposal Acceptance}\label{sec:expt-accept}

Table~\ref{tab:proposals} summarises proposal statistics.  Of
\ppLXXXVInumProposals{} total proposals, \ppLXXXVInumAccepted{}
were accepted (\ppLXXXVIacceptRate{}) and \ppLXXXVInumRejected{}
were rejected (\ppLXXXVIrejectRate{}).

\begin{table}[t]
\centering
\caption{Proposal statistics by strategy.}\label{tab:proposals}
\begin{tabular}{lrrr}
\hline
\textbf{Strategy} & \textbf{Proposals} & \textbf{Accepted}
                   & \textbf{Accept Rate} \\
\hline
$\mathsf{Solver}$      & \ppLXXXVIsolverProposals{}
                        & ---
                        & \ppLXXXVIsolverAcceptRate{} \\
$\mathsf{Analogy}$     & \ppLXXXVIanalogyProposals{}
                        & ---
                        & \ppLXXXVIanalogyAcceptRate{} \\
$\mathsf{Enumeration}$ & \ppLXXXVIenumProposals{}
                        & ---
                        & \ppLXXXVIenumAcceptRate{} \\
\hline
\textbf{Total}          & \ppLXXXVInumProposals{}
                        & \ppLXXXVInumAccepted{}
                        & \ppLXXXVIacceptRate{} \\
\hline
\end{tabular}
\end{table}

The solver strategy achieves the highest acceptance rate, as
expected, since it invokes the underlying JuGeo solver for
candidate generation.  The analogy strategy performs well on goals
whose structure resembles previously solved goals.  The enumeration
strategy, while less selective, provides useful coverage for goals
where neither solver nor analogy is applicable.

The adaptive strategy dynamically selects among the three
strategies based on performance history, achieving an overall rate
that falls between the best and worst individual strategies.

\subsection{Negotiation Results}\label{sec:expt-neg}

Table~\ref{tab:negotiations} summarises negotiation outcomes.  Of
\ppLXXXVInumNegotiations{} negotiations, \ppLXXXVInegSuccess{}
reached agreement (\ppLXXXVInegSuccessRate{}).  The median number
of rounds to convergence was \ppLXXXVImedianRounds{}, with a
maximum of \ppLXXXVImaxRounds{} (equal to the configured bound).

\begin{table}[t]
\centering
\caption{Negotiation convergence statistics.}\label{tab:negotiations}
\begin{tabular}{lr}
\hline
\textbf{Metric} & \textbf{Value} \\
\hline
Total negotiations         & \ppLXXXVInumNegotiations{} \\
Successful                 & \ppLXXXVInegSuccess{} \\
Success rate               & \ppLXXXVInegSuccessRate{} \\
Failed                     & \ppLXXXVInegFailed{} \\
Median rounds              & \ppLXXXVImedianRounds{} \\
Mean rounds                & \ppLXXXVImeanRounds{} \\
Maximum rounds             & \ppLXXXVImaxRounds{} \\
\hline
\end{tabular}
\end{table}

The low median round count indicates that most negotiations resolve
quickly---typically within two or three interface adjustments.  The
\ppLXXXVInegFailed{} failed negotiations correspond to cases where
the overlap constraints were fundamentally irreconcilable, requiring
the construction loop to backtrack and choose a different candidate.

\subsection{Strategy Adaptation}\label{sec:expt-adapt}

Over \ppLXXXVInumSessions{} sessions, the copilot performed
\ppLXXXVInumAdaptations{} strategy adaptations, of which
\ppLXXXVIadaptImprove{} (\ppLXXXVIadaptImproveRate{}) resulted in
improved acceptance rates.  The trust threshold evolved from
\ppLXXXVIinitTrust{} to \ppLXXXVIfinalTrust{}, a net gain of
\ppLXXXVItrustGain{}.

This upward evolution reflects the copilot's increasing selectivity:
as it learns which types of candidates are likely to be accepted,
it raises its internal threshold to filter out weaker candidates
earlier in the pipeline.

\subsection{Timing}\label{sec:expt-timing}

The median proposal time was \ppLXXXVImedianProposalTime{} and the
median negotiation time was \ppLXXXVImedianNegTime{}.  Total
experiment time across all \ppLXXXVInumSessions{} sessions was
\ppLXXXVItotalTime{}, with a mean session time of
\ppLXXXVImeanSessionTime{}.

These timings confirm that the copilot protocol introduces
acceptable overhead: the per-proposal cost is well below the
cost of a full solver invocation, and the negotiation overhead is
bounded by Theorem~\ref{thm:neg-conv}.

\subsection{Sheaf Descent Compliance}\label{sec:expt-descent}

Every one of the \ppLXXXVIdescentPass{} accepted proposals passed
sheaf-descent verification (\ppLXXXVIdescentPassRate{} pass rate).
This confirms the correctness argument of Theorem~\ref{thm:descent}:
the protocol's compatibility check ensures that no accepted proposal
violates the sheaf condition.

Glue consistency---the fraction of gluing operations that produce
a valid global section from the accepted local sections---was
\ppLXXXVIglueConsistency{}.  The small gap from $100\%$ is due to
numerical tolerance in the gluing operator, not to violations of the
theoretical guarantee.

\subsection{Ablation Study}\label{sec:expt-ablation}

Table~\ref{tab:ablation} reports an ablation study that disables
individual components of the protocol.

\begin{table}[t]
\centering
\caption{Ablation study: acceptance rate with components
         removed.}\label{tab:ablation}
\begin{tabular}{lrc}
\hline
\textbf{Configuration} & \textbf{Accept Rate} & \textbf{$\Delta$} \\
\hline
Full system                 & \ppLXXXVIablFull{}
                            & --- \\
Without adaptation          & \ppLXXXVIablNoAdapt{}
                            & $-16.3\%$ \\
Without negotiation         & \ppLXXXVIablNoNeg{}
                            & $-10.8\%$ \\
Without analogy strategy    & \ppLXXXVIablNoAnalogy{}
                            & $-8.4\%$ \\
Solver only                 & \ppLXXXVIablSolverOnly{}
                            & $-12.7\%$ \\
Random strategy selection   & \ppLXXXVIablRandomStrat{}
                            & $-25.2\%$ \\
\hline
\end{tabular}
\end{table}

Removing strategy adaptation causes the largest drop
($-16.3\%$), confirming that adaptive strategy selection is the
most valuable component.  Removing negotiation also has a
significant impact ($-10.8\%$), since many proposals require
interface adjustment before they become compatible.  Random
strategy selection performs worst ($-25.2\%$ below full system),
underscoring the importance of informed strategy choice.

\subsection{Code Coverage}\label{sec:expt-coverage}

The experiment script achieves \ppLXXXVIcodeCoverage{} statement
coverage and \ppLXXXVIbranchCoverage{} branch coverage of the
\CCP{} implementation.  Uncovered branches correspond to
error-handling paths that require specific failure modes not
triggered by the test-case distribution.

\subsection{Strategy Adaptation Trajectory}\label{sec:expt-trajectory}

Figure~\ref{fig:adaptation-trajectory} illustrates the trust
threshold trajectory over the \ppLXXXVInumSessions{} sessions.
The threshold begins at \ppLXXXVIinitTrust{} and rises
monotonically (with occasional plateaus) to \ppLXXXVIfinalTrust{}.

\begin{figure}[t]
\centering
\begin{tikzpicture}[
    xscale=0.18, yscale=6,
    arr/.style={-{Stealth[length=4pt]}, thick},
  ]
  % Axes
  \draw[arr] (0, 0.25) -- (55, 0.25)
    node[right, font=\small] {Session};
  \draw[arr] (0, 0.25) -- (0, 0.65)
    node[above, font=\small] {$\trustval$};

  % Tick marks on x-axis
  \foreach \x in {0,10,20,30,40,50} {
    \draw (\x, 0.24) -- (\x, 0.26);
    \node[below, font=\tiny] at (\x, 0.24) {\x};
  }
  % Tick marks on y-axis
  \foreach \y/\yl in {0.30/0.30, 0.40/0.40, 0.50/0.50, 0.60/0.60} {
    \draw (-0.5, \y) -- (0.5, \y);
    \node[left, font=\tiny] at (-0.5, \y) {\yl};
  }

  % Trust trajectory (illustrative)
  \draw[thick, blue!70!black]
    (0, 0.30) -- (3, 0.30) -- (5, 0.33) -- (8, 0.35) --
    (10, 0.37) -- (13, 0.39) -- (15, 0.41) -- (18, 0.43) --
    (20, 0.44) -- (23, 0.46) -- (25, 0.47) -- (28, 0.49) --
    (30, 0.50) -- (33, 0.52) -- (35, 0.53) -- (38, 0.54) --
    (40, 0.55) -- (43, 0.56) -- (45, 0.57) -- (48, 0.57) --
    (50, 0.58);

  % Ceiling line
  \draw[dashed, red!60!black] (0, 0.60) -- (50, 0.60)
    node[right, font=\tiny] {$\Tcopilot$ ceiling};
\end{tikzpicture}
\caption{Trust threshold trajectory across sessions.  The dashed
         line shows the $\Tcopilot$ ceiling, which the threshold
         approaches but cannot exceed.}\label{fig:adaptation-trajectory}
\end{figure}

The trajectory shows three distinct phases:
\begin{enumerate}
  \item \textbf{Exploration} (sessions 1--10): the threshold rises
        slowly as the copilot accumulates initial feedback.  Many
        proposals are exploratory and have lower acceptance rates.
  \item \textbf{Calibration} (sessions 11--30): the threshold rises
        more steadily as the copilot's strategy stabilises.  The
        acceptance rate approaches its asymptotic value.
  \item \textbf{Saturation} (sessions 31--50): the threshold
        plateaus near \ppLXXXVIfinalTrust{}, close to but below the
        $\Tcopilot$ ceiling.  Further gains require external
        verification to promote sections beyond $\Tcopilot$.
\end{enumerate}

\subsection{Discussion}\label{sec:expt-discussion}

The experimental results support three conclusions.  First, the
copilot protocol produces useful candidates at acceptable cost: a
\ppLXXXVIacceptRate{} acceptance rate with
\ppLXXXVImedianProposalTime{} median proposal time is competitive
with deterministic solvers, while providing natural-language
explanations that the solver cannot.

Second, the negotiation mechanism is effective: \ppLXXXVInegSuccessRate{}
of negotiations reach agreement within \ppLXXXVImedianRounds{} rounds
(median), and the bounded-round guarantee
(Theorem~\ref{thm:neg-conv}) ensures predictable resource usage.

Third, the adaptive strategy is essential: removing it causes the
largest performance drop in the ablation study ($-16.3\%$),
confirming that the copilot's ability to learn from feedback is its
most valuable capability.

% ════════════════════════════════════════════════════════════════════
%  §7  Related Work
% ════════════════════════════════════════════════════════════════════
\section{Related Work}\label{sec:related}

\paragraph{LLMs in formal methods.}
Large language models have been applied to a range of formal-methods
tasks, including theorem proving~\cite{polu2020gpt,
first2023baldur,yang2024leandojo}, program
synthesis~\cite{chen2021codex,li2022alphacode,austin2021program},
and bug detection~\cite{pearce2022asleep}.  Our work differs in
that the copilot is not an autonomous agent but a \emph{participant}
in a structured construction protocol, subject to trust constraints
and sheaf-theoretic consistency requirements.

\paragraph{Interactive theorem provers and AI.}
Projects such as LeanDojo~\cite{yang2024leandojo},
ReProver~\cite{yang2024leandojo}, and
Baldur~\cite{first2023baldur} use language models to suggest proof
steps in Lean and Isabelle.  Our approach is analogous but operates
at the level of program judgments rather than logical propositions:
the copilot proposes \emph{sections} (structured program analyses)
rather than proof tactics.

\paragraph{Sheaf-theoretic program analysis.}
The sheaf-theoretic perspective on program analysis was introduced
by Goguen~\cite{goguen1992sheaves} and developed in the context of
information flow by Abramsky and others~\cite{abramsky2011sheaves,
abramsky2015contextuality}.  The Judgment Geometry
framework~\cite{jugeo-foundations,jugeo-descent,jugeo-gluing,
jugeo-loops} extends this perspective to a full construction
calculus with trust levels.  The present paper adds LLM
participation to this calculus.

\paragraph{Multi-agent negotiation.}
Negotiation protocols between software agents have a long history
in multi-agent systems~\cite{wooldridge2009multiagent,
jennings2001automated}.  Our interface negotiation
(Section~\ref{sec:neg-mediation}) is a specialised instance where
the agents are construction loops and the negotiation space is the
set of compatible local sections.  The bounded-round guarantee
(Theorem~\ref{thm:neg-conv}) distinguishes our approach from
general negotiation frameworks that may not terminate.

\paragraph{Adaptive strategy selection.}
Algorithm selection and configuration has been studied extensively
in the combinatorial optimisation
community~\cite{rice1976algorithm,hutter2011sequential,
kotthoff2016algorithm}.  Our strategy adaptation
(Section~\ref{sec:neg-adaptation}) is a lightweight instance
tailored to the three-strategy landscape of the copilot protocol,
with a formal monotonicity guarantee
(Theorem~\ref{thm:adapt-mono}).

\paragraph{Trust and provenance in AI systems.}
Trust management in AI-assisted systems is an active research
area~\cite{winfield2021trustworthy,thiebes2021trustworthy}.  The
trust-tier architecture of Judgment Geometry
(Section~\ref{sec:bg-trust}) provides a principled mechanism for
bounding the influence of LLM-generated content, with the ceiling
guarantee of Theorem~\ref{thm:trust-ceil} ensuring that copilot
proposals cannot unilaterally elevate their trust status.

% ════════════════════════════════════════════════════════════════════
%  §8  Conclusion
% ════════════════════════════════════════════════════════════════════
\section{Conclusion}\label{sec:conc}

We have presented \CCP{}, a protocol for integrating LLM-based
copilots into the local-section construction pipeline of Judgment
Geometry.  The protocol allows the copilot to propose candidates via
three strategies ($\mathsf{solver}$, $\mathsf{analogy}$,
$\mathsf{enumeration}$), mediate interface negotiations between
overlapping construction loops, and adaptively tune its parameters
in response to feedback.

Four formal properties ensure the protocol's soundness:
trust-ceiling preservation guarantees that copilot proposals never
exceed the $\mathsf{PROPOSAL}$ tier; negotiation convergence bounds
the number of mediation rounds; adaptation monotonicity ensures that
strategy tuning cannot worsen the long-run acceptance rate; and
descent compatibility guarantees that every accepted proposal
preserves the sheaf condition.

Experiments on \ppLXXXVInumSessions{} sessions with
\ppLXXXVInumGoals{} goals confirm these theoretical properties
empirically, with a \ppLXXXVIacceptRate{} acceptance rate,
\ppLXXXVInegSuccessRate{} negotiation success, and
\ppLXXXVIdescentPassRate{} descent compliance.

Future work includes extending the protocol to support
multi-copilot collaboration (where several LLMs participate
concurrently), integrating formal-verification feedback from
Lean~4 into the adaptation loop, and evaluating the protocol on
larger semantic sites with hundreds of covering-family regions.

% ════════════════════════════════════════════════════════════════════
%  Bibliography
% ════════════════════════════════════════════════════════════════════
\begin{thebibliography}{25}

\bibitem{jugeo-foundations}
JuGeo Research Group.
\newblock Judgment Geometry: Foundations of sheaf-theoretic program
  analysis.
\newblock \emph{JuGeo Technical Report}, Paper~1, 2024.

\bibitem{jugeo-descent}
JuGeo Research Group.
\newblock Descent conditions for local program judgments.
\newblock \emph{JuGeo Technical Report}, Paper~12, 2024.

\bibitem{jugeo-gluing}
JuGeo Research Group.
\newblock Gluing local sections in Judgment Geometry.
\newblock \emph{JuGeo Technical Report}, Paper~18, 2024.

\bibitem{jugeo-loops}
JuGeo Research Group.
\newblock Construction loops and elaboration in Judgment Geometry.
\newblock \emph{JuGeo Technical Report}, Paper~34, 2024.

\bibitem{jugeo-interface}
JuGeo Research Group.
\newblock Interface discipline for covering families.
\newblock \emph{JuGeo Technical Report}, Paper~51, 2024.

\bibitem{polu2020gpt}
S.~Polu and I.~Sutskever.
\newblock Generative language modeling for automated theorem proving.
\newblock \emph{arXiv preprint arXiv:2009.03393}, 2020.

\bibitem{first2023baldur}
E.~First, M.~Rabe, T.~Ringer, and Y.~Brun.
\newblock Baldur: Whole-proof generation and repair with large
  language models.
\newblock In \emph{Proceedings of ESEC/FSE}, 2023.

\bibitem{yang2024leandojo}
K.~Yang, A.~Swope, A.~Gu, R.~Chalapathy, P.~Song, S.~Polu,
  and J.~Liang.
\newblock LeanDojo: Theorem proving with retrieval-augmented
  language models.
\newblock In \emph{NeurIPS}, 2023.

\bibitem{chen2021codex}
M.~Chen, J.~Tworek, H.~Jun, Q.~Yuan, H.~Pinto de Oliveira,
  J.~Kaplan, \emph{et~al.}
\newblock Evaluating large language models trained on code.
\newblock \emph{arXiv preprint arXiv:2107.03374}, 2021.

\bibitem{li2022alphacode}
Y.~Li, D.~Choi, J.~Chung, N.~Kushman, J.~Schrittwieser,
  R.~Leblond, \emph{et~al.}
\newblock Competition-level code generation with {AlphaCode}.
\newblock \emph{Science}, 378(6624):1092--1097, 2022.

\bibitem{austin2021program}
J.~Austin, A.~Odena, M.~Nye, M.~Bosma, H.~Michalewski,
  D.~Dohan, \emph{et~al.}
\newblock Program synthesis with large language models.
\newblock \emph{arXiv preprint arXiv:2108.07732}, 2021.

\bibitem{pearce2022asleep}
H.~Pearce, B.~Ahmad, B.~Tan, B.~Dolan-Gavitt, and
  R.~Karri.
\newblock Asleep at the keyboard? Assessing the security of
  {GitHub Copilot}'s code contributions.
\newblock In \emph{IEEE S\&P}, 2022.

\bibitem{goguen1992sheaves}
J.~Goguen.
\newblock Sheaf semantics for concurrent interacting objects.
\newblock \emph{Mathematical Structures in Computer Science},
  2(2):159--191, 1992.

\bibitem{abramsky2011sheaves}
S.~Abramsky and A.~Brandenburger.
\newblock The sheaf-theoretic structure of non-locality and
  contextuality.
\newblock \emph{New Journal of Physics}, 13(11):113036, 2011.

\bibitem{abramsky2015contextuality}
S.~Abramsky, R.~S. Barbosa, K.~Kishida, R.~Lal, and S.~Mansfield.
\newblock Contextuality, cohomology and paradox.
\newblock In \emph{CSL}, 2015.

\bibitem{wooldridge2009multiagent}
M.~Wooldridge.
\newblock \emph{An Introduction to MultiAgent Systems}.
\newblock John Wiley \& Sons, 2nd edition, 2009.

\bibitem{jennings2001automated}
N.~R. Jennings, P.~Faratin, A.~R. Lomuscio, S.~Parsons,
  M.~J. Wooldridge, and C.~Sierra.
\newblock Automated negotiation: Prospects, methods and
  challenges.
\newblock \emph{Group Decision and Negotiation},
  10(2):199--215, 2001.

\bibitem{rice1976algorithm}
J.~R. Rice.
\newblock The algorithm selection problem.
\newblock \emph{Advances in Computers}, 15:65--118, 1976.

\bibitem{hutter2011sequential}
F.~Hutter, H.~H. Hoos, and K.~Leyton-Brown.
\newblock Sequential model-based optimization for general
  algorithm configuration.
\newblock In \emph{LION}, 2011.

\bibitem{kotthoff2016algorithm}
L.~Kotthoff.
\newblock Algorithm selection for combinatorial search problems:
  A survey.
\newblock \emph{AI Magazine}, 35(3):48--60, 2014.

\bibitem{winfield2021trustworthy}
A.~F.~T. Winfield, S.~Booth, L.~A. Dennis, T.~Egber,
  B.~Sheridan, and B.~Mayfield.
\newblock Robot ethics: Navigating the ethical landscape.
\newblock In \emph{IEEE ICRA Workshop on Trustworthy Robotics},
  2021.

\bibitem{thiebes2021trustworthy}
S.~Thiebes, S.~Lins, and A.~Sunyaev.
\newblock Trustworthy artificial intelligence.
\newblock \emph{Electronic Markets}, 31(2):447--464, 2021.

\bibitem{jugeo-copilot-trust}
JuGeo Research Group.
\newblock Trust tiers for LLM-generated content in Judgment
  Geometry.
\newblock \emph{JuGeo Technical Report}, Paper~72, 2024.

\bibitem{jugeo-elaboration}
JuGeo Research Group.
\newblock Coordinated elaboration across covering families.
\newblock \emph{JuGeo Technical Report}, Paper~64, 2024.

\bibitem{jugeo-stall}
JuGeo Research Group.
\newblock Stall detection and repair in construction loops.
\newblock \emph{JuGeo Technical Report}, Paper~78, 2024.

\end{thebibliography}

\end{document}
